\chapter{Hinweise zur Integration des Backends}\label{chapter:backend-integration}

Die Details dazu, wo und wie das Backend integriert werden soll, wurden bereits in den vorherigen Kapiteln zu den einzelnen Seiten der Webapp erläutert. Hier folgt nun noch eine kurze Zusammenfassung:\\

Generell wurde an jeder relevanten Stelle im Code ein ``TODO BACKEND'' Kommentar hinterlassen, oft schon mit sehr konkreten Vorschlägen wie die Funktion des Backends heißen könnte, welche Daten übergeben müssen und wie diese weiterverarbeitet werden müssen. Auch mit Backend dann nicht mehr benötigte Teile des Codes sind so gekennzeichnet. Die meisten Änderungen sind in der \texttt{TrackingEditor} Komponente nötig bzw. beschrieben. Dort gibt es einige state Variablen, die zukünftig vom Backend verwaltet und abgefragt werden sollen, diese sind mit einem ``BACKEND'' Kommentar markiert. Außerdem sind generell alle JSON Dateien im \texttt{/src/data/} Ordner im Prinzip manuell erstellte (oder im Fall der Trackingdaten zur Verfügung gestellte reale Daten) Antworten des Backends. Diese zeigen also die genaue Struktur, die die Antworten des Backends haben müssen, damit diese direkt ohne weitere Änderungen von der Webapp weiter verarbeitet und angezeigt werden können. Außerdem müssen auch die im gleichen Ordner abgelegten JPG Bilder, also der Thumbnail der einzelnen Videos für die \texttt{FileOverview} Komponente sowie die Videoframes, die in \texttt{TrackingEditor} angezeigt werden, vom Backend gesendet werden, jeweils zusammen mit den anderen in den jeweiligen Komponenten verwendeten Daten.\\
\\
Bei der Entwicklung der Analyse Seite wurde Wert darauf gelegt, den Code möglichst allgemein zu halten, damit einfach weitere Analysen ohne große Änderungen im Frontend angezeigt werden können. Deshalb werden z.B. die Beschriftungen der einzelnen Achsen für die Diagramme direkt vom Backend übergeben. So können einfach durch das Übergeben anderer Daten und Beschriftungen beliebige Balkendiagramme angezeigt werden. Die Auswahl, welche Diagramme angezeigt werden sollen muss dann allerdings noch ergänzt werden. \\
\\
Ein eventuelles Streamen des Videos konnte in Ermangelung eines Backends ebenfalls nicht getestet werden. Meiner Einschätzung nach könnte es schwierig sein, die angezeigten Trackingdaten mit den aktuellen Videoframes konsistent zu halten, bzw. immer abzufragen welcher Videoframe aktuell angezeigt wird. Eine Möglichkeit, dieses Problem zu umgehen, wäre die Trackingdaten im Backend direkt mit in ein Video zu rendern, sodass einfach nur ein Video angezeigt wird. Dieser Renderprozess müsste evtl. vor dem Streamen des Videos bereits beendet sein, um ein Streamen in Echtzeit zu ermöglichen. Sobald das Video angehalten wird, könnte dann wieder der einzelne Frame ohne die Trackingdaten gesendet werden, sodass die Bounding Boxes und Ecken wieder von der Webapp selbst angezeigt werden und somit bearbeitbar sind. \\
Alternativ oder zusätzlich könnte auch die Framerate des Videos oder auch nur der Trackingdaten reduziert werden. Möglicherweise ist aber auch ein Streamen des Videos und ein Anzeigen der Trackingdaten ohne größere Probleme möglich. Dann wäre allerdings noch die Frage zu klären, wie das Backend mit mehreren laufenden Instanzen der Webapp zurecht kommt bzw. wie sehr es dafür skaliert werden müsste.
